\begin{abcliste}
\abc Split the velocity of the particle into a component along the field lines and a component across them. Across the field lines, the magnetic force (perpendicular to the velocity and to the field) makes the particle move on a circle. Along the field lines there is no force, so this component of the velocity stays the same. Together, the particles move on $\result{\text{helices (spirals) along the field lines}}$: not straight along the field lines, and not on circles around the Earth's centre.

\abc The particles follow the field lines, and the field lines of the Earth lead to the regions around the magnetic poles, where they enter the atmosphere. $\result{\text{The field lines lead the particles there}}$, and there the air glows. (It has nothing to do with the colder air or with the sunlight at the poles.)

\abc Towards the poles the field lines crowd together: the field gets stronger and the circles get tighter. Most particles are turned back before they reach the ground (a \emph{magnetic mirror}) and bounce back and forth between north and south. $\result{\text{The stronger field near the poles turns them back.}}$

Note: the magnetic force does no work, since it is always perpendicular to the velocity. It cannot slow the particles down; it only changes the direction of their motion.
\end{abcliste}
